\begingroup
\let\cleardoublepage\relax


\chapter*{Grundlegende Methoden in Java}
\label{ch:grundlegende-methoden-in-java}
\addcontentsline{toc}{chapter}{Grundlegende Methoden in Java}
\endgroup


\section*{String-Methoden}
\addcontentsline{toc}{section}{String-Methoden}

\subsection*{length() - Länge eines Strings}

\subsubsection*{Was macht die Methode?}\\
Gibt die Anzahl der Zeichen im String zurück.

\subsubsection*{Syntax:}
\begin{lstlisting}
int laenge = text.length();
\end{lstlisting}

\subsubsection*{Beispiele:}
\begin{lstlisting}
String name = "Anna";
int laenge = name.length();  // laenge = 4

String wort = "Hallo Welt";
int laenge2 = wort.length();  // laenge2 = 10 (Leerzeichen zaehlt mit!)

String leer = "";
int laenge3 = leer.length();  // laenge3 = 0

String satz = "Java";
for (int i = 0; i < satz.length(); i++) {
    // Schleife laeuft 4 mal (von 0 bis 3)
}
\end{lstlisting}

\begin{achtung}
	\begin{itemize}
		\item \texttt{length()} zählt alle Zeichen: Der Text "Hi" hat die Länge 2
		\item Leerzeichen und Sonderzeichen zählen mit
		\item Bei Strings: \texttt{.length()} mit Klammern
		\item Bei Arrays: \texttt{.length} ohne Klammern - das ist ein wichtiger Unterschied!
	\end{itemize}
\end{achtung}

% ============================================================
\newpage


\subsection*{charAt(int index) - Zeichen an Position}

\subsubsection*{Was macht die Methode?}\\
Gibt dir das Zeichen an einer bestimmten Stelle im Text zurück. Wichtig: Die Zählung beginnt bei 0!

\subsubsection*{Syntax:}
\begin{lstlisting}
char zeichen = text.charAt(index);
\end{lstlisting}

\subsubsection*{Beispiele:}
\begin{lstlisting}
String wort = "Hallo";
char erstes = wort.charAt(0);   // erstes = 'H'
char zweites = wort.charAt(1);  // zweites = 'a'
char letztes = wort.charAt(4);  // letztes = 'o'

// Alle Zeichen durchlaufen
String text = "Java";
for (int i = 0; i < text.length(); i++) {
    char c = text.charAt(i);
    System.out.println(c);  // Gibt aus: J, a, v, a
}

// Pruefung einzelner Zeichen
String eingabe = "A1B2";
if (eingabe.charAt(0) == 'A') {
    System.out.println("Beginnt mit A");
}

// Zeichen vergleichen
String code = "12345";
if (code.charAt(2) == '3') {
    System.out.println("An Position 2 steht eine 3");
}
\end{lstlisting}

\begin{achtung}
	\begin{itemize}
		\item Die Nummerierung beginnt bei 0
		\item Die Nummer muss zwischen 0 und \texttt{text.length() - 1} liegen
		\item Wenn die Nummer zu groß oder zu klein ist, stürzt das Programm ab
		\item Du bekommst ein einzelnes Zeichen (char) zurück, keinen Text (String)
	\end{itemize}
\end{achtung}

% ============================================================
\newpage


\section*{Ausgabe auf der Konsole}
\addcontentsline{toc}{section}{Ausgabe auf der Konsole}

Mit \texttt{System.out.println()} und \texttt{System.out.print()} kannst du Text und Zahlen auf der Konsole ausgeben.


\subsection*{println() - Ausgabe mit Zeilenumbruch}
\addcontentsline{toc}{subsection}{println() - Ausgabe mit Zeilenumbruch}

\subsubsection*{Was macht die Methode?}\\
Gibt Text, Zahlen oder andere Werte auf der Konsole aus und springt danach in die nächste Zeile.

\subsubsection*{Syntax:}
\begin{lstlisting}
System.out.println(wert);
\end{lstlisting}

\subsubsection*{Beispiele:}
\begin{lstlisting}
// Text ausgeben
System.out.println("Hallo Welt!");
// Ausgabe: Hallo Welt!

// Zahlen ausgeben
System.out.println(42);
// Ausgabe: 42

System.out.println(3.14);
// Ausgabe: 3.14

// Variablen ausgeben
String name = "Anna";
System.out.println(name);
// Ausgabe: Anna

int alter = 25;
System.out.println(alter);
// Ausgabe: 25

// Text und Variablen kombinieren
System.out.println("Ich bin " + alter + " Jahre alt");
// Ausgabe: Ich bin 25 Jahre alt

System.out.println("Hallo " + name + "!");
// Ausgabe: Hallo Anna!

// Mehrere Ausgaben hintereinander
System.out.println("Erste Zeile");
System.out.println("Zweite Zeile");
System.out.println("Dritte Zeile");
// Ausgabe:
// Erste Zeile
// Zweite Zeile
// Dritte Zeile

// Berechnungen direkt ausgeben
System.out.println(5 + 3);
// Ausgabe: 8

System.out.println("Ergebnis: " + (10 * 2));
// Ausgabe: Ergebnis: 20

// Leere Zeile ausgeben
System.out.println();
// Ausgabe: (eine leere Zeile)

// Boolean-Werte ausgeben
boolean istWahr = true;
System.out.println(istWahr);
// Ausgabe: true
\end{lstlisting}

\begin{achtung}
	\begin{itemize}
		\item Nach der Ausgabe springt der Cursor automatisch in die nächste Zeile
		\item Du kannst Text, Zahlen, Variablen und sogar Berechnungen ausgeben
		\item Zum Verbinden von Text und Variablen benutzt du das Plus-Zeichen (+)
		\item Ohne Wert macht \texttt{println()} einfach eine neue Zeile
	\end{itemize}
\end{achtung}

% ============================================================
\newpage


\subsection*{print() - Ausgabe ohne Zeilenumbruch}
\addcontentsline{toc}{subsection}{print() - Ausgabe ohne Zeilenumbruch}

\subsubsection*{Was macht die Methode?}\\
Gibt Text, Zahlen oder andere Werte auf der Konsole aus, bleibt aber in der gleichen Zeile.

\subsubsection*{Syntax:}
\begin{lstlisting}
System.out.print(wert);
\end{lstlisting}

\subsubsection*{Beispiele:}
\begin{lstlisting}
// Text ausgeben ohne Zeilenumbruch
System.out.print("Hallo ");
System.out.print("Welt!");
// Ausgabe: Hallo Welt!

// Zahlen ausgeben
System.out.print(1);
System.out.print(2);
System.out.print(3);
// Ausgabe: 123

// Mit Leerzeichen trennen
System.out.print("Eins ");
System.out.print("Zwei ");
System.out.print("Drei");
// Ausgabe: Eins Zwei Drei

// Zähler in einer Zeile
for (int i = 1; i <= 5; i++) {
    System.out.print(i + " ");
}
// Ausgabe: 1 2 3 4 5

// Sterne in einer Reihe ausgeben
for (int i = 0; i < 10; i++) {
    System.out.print("*");
}
// Ausgabe: **********

// Kombination mit println()
System.out.print("Name: ");
System.out.println("Anna");
System.out.print("Alter: ");
System.out.println(25);
// Ausgabe:
// Name: Anna
// Alter: 25
\end{lstlisting}

\begin{achtung}
	\begin{itemize}
		\item Der Cursor bleibt in der gleichen Zeile - es gibt keinen Zeilenumbruch
		\item Nützlich, wenn du mehrere Dinge in eine Zeile schreiben willst
		\item Oft wird \texttt{print()} mit \texttt{println()} kombiniert
		\item Denk daran, Leerzeichen einzufügen, wenn du Wörter trennen willst
	\end{itemize}
\end{achtung}

% ============================================================
\newpage


\subsection*{Der Unterschied: print() vs. println()}

\subsubsection*{Direkter Vergleich:}
\begin{lstlisting}
// Mit println() - jedes Mal eine neue Zeile
System.out.println("A");
System.out.println("B");
System.out.println("C");
// Ausgabe:
// A
// B
// C

// Mit print() - alles in einer Zeile
System.out.print("A");
System.out.print("B");
System.out.print("C");
// Ausgabe:
// ABC

// Gemischt - Kombination aus beiden
System.out.print("Vorname: ");
System.out.println("Max");
System.out.print("Nachname: ");
System.out.println("Mustermann");
// Ausgabe:
// Vorname: Max
// Nachname: Mustermann
\end{lstlisting}

\subsubsection*{Wann benutzt du was?}

\begin{itemize}
	\item \textbf{println()}: Wenn jede Ausgabe in einer eigenen Zeile stehen soll
	\item \textbf{print()}: Wenn du mehrere Dinge in einer Zeile ausgeben willst
	\item \textbf{Kombination}: print() für die Beschriftung, println() für den Wert
\end{itemize}

\subsubsection*{Praktisches Beispiel:}
\begin{lstlisting}
// Quadratzahlen ausgeben
System.out.println("Quadratzahlen von 1 bis 5:");
for (int i = 1; i <= 5; i++) {
    System.out.print(i + "² = ");
    System.out.println(i * i);
}
// Ausgabe:
// Quadratzahlen von 1 bis 5:
// 1² = 1
// 2² = 4
// 3² = 9
// 4² = 16
// 5² = 25

// Muster ausgeben
for (int zeile = 1; zeile <= 3; zeile++) {
    for (int spalte = 1; spalte <= 5; spalte++) {
        System.out.print("* ");
    }
    System.out.println();  // Neue Zeile nach jeder Reihe
}
// Ausgabe:
// * * * * *
// * * * * *
// * * * * *
\end{lstlisting}

% ============================================================
\newpage


\section*{ConsoleReader - Eingaben von der Konsole}
\addcontentsline{toc}{section}{ConsoleReader - Eingaben von der Konsole}

Der \texttt{ConsoleReader} ist ein Werkzeug, mit dem du Eingaben vom Benutzer einlesen kannst.
Um ihn zu benutzen, schreibst du am Anfang deines Programms:

\begin{lstlisting}
import static prog.ConsoleReader.*;
\end{lstlisting}

Alle Eingabe-Methoden gibt es in zwei Varianten:
\begin{itemize}
	\item Ohne Text: Wartet einfach auf die Eingabe
	\item Mit Text: Zeigt zuerst eine Aufforderung an (z.B. "Wie heißt du?")
\end{itemize}

% ============================================================
\newpage


\subsection*{readString() - String einlesen}
\addcontentsline{toc}{subsection}{readString() - String einlesen}

\subsubsection*{Was macht die Methode?}\\
Liest eine Zeile Text von der Konsole ein und gibt sie als String zurück.

\subsubsection*{Syntax:}
\begin{lstlisting}
String text = readString();                // Ohne Prompt
String text = readString("Eingabetext");   // Mit Prompt
\end{lstlisting}

\subsubsection*{Beispiele:}
\begin{lstlisting}
// Einfache Eingabe
String name = readString("Wie heisst du? ");
System.out.println("Hallo " + name);

// Mehrere Eingaben
String vorname = readString("Vorname");
String nachname = readString("Nachname");
String vollername = vorname + " " + nachname;

// Ohne Prompt
System.out.println("Gib einen Text ein:");
String eingabe = readString();

// In Schleife verwenden
String antwort;
do {
    antwort = readString("Weiter? (ja/nein)");
} while (!antwort.equals("nein"));
\end{lstlisting}

\begin{achtung}
	\begin{itemize}
		\item Liest alles, was der Benutzer eingibt, bis er Enter drückt
		\item Du bekommst immer Text (String) zurück, auch wenn jemand "123" eingibt
		\item Wenn du Zahlen brauchst, benutze \texttt{readInt()} oder \texttt{readDouble()}
		\item Wenn der Benutzer nur Enter drückt, bekommst du einen leeren Text \texttt{""} zurück
	\end{itemize}
\end{achtung}

% ============================================================
\newpage


\subsection*{readInt() - Ganzzahl einlesen}
\addcontentsline{toc}{subsection}{readInt() - Ganzzahl einlesen}

\subsubsection*{Was macht die Methode?}\\
Liest eine Ganzzahl von der Konsole ein und gibt sie als \texttt{int} zurück.

\subsubsection*{Syntax:}
\begin{lstlisting}
int zahl = readInt();                      // Ohne Prompt
int zahl = readInt("Gib eine Zahl ein");   // Mit Prompt
\end{lstlisting}

\subsubsection*{Beispiele:}
\begin{lstlisting}
// Alter einlesen
int alter = readInt("Wie alt bist du? ");
System.out.println("Du bist " + alter + " Jahre alt");

// Berechnungen
int a = readInt("Erste Zahl");
int b = readInt("Zweite Zahl");
int summe = a + b;
System.out.println("Summe: " + summe);

// Array-Groesse einlesen
int anzahl = readInt("Wie viele Zahlen? ");
int[] zahlen = new int[anzahl];
for (int i = 0; i < zahlen.length; i++) {
    zahlen[i] = readInt("Zahl " + (i + 1));
}

// Schleifensteuerung
int wiederholungen = readInt("Wie oft wiederholen? ");
for (int i = 0; i < wiederholungen; i++) {
    System.out.println("Durchlauf " + (i + 1));
}
\end{lstlisting}

\begin{achtung}
	\begin{itemize}
		\item Wenn jemand Text statt einer Zahl eingibt, stürzt das Programm ab
		\item Funktioniert nur für ganze Zahlen von etwa -2 Milliarden bis +2 Milliarden
		\item Für größere Zahlen: \texttt{readLong()} benutzen
		\item Für Kommazahlen: \texttt{readDouble()} benutzen
	\end{itemize}
\end{achtung}

% ============================================================
\newpage


\subsection*{readDouble() - Kommazahl einlesen}
\addcontentsline{toc}{subsection}{readDouble() - Kommazahl einlesen}

\subsubsection*{Was macht die Methode?}\\
Liest eine Kommazahl von der Konsole ein und gibt sie als \texttt{double} zurück.

\subsubsection*{Syntax:}
\begin{lstlisting}
double zahl = readDouble();                    // Ohne Prompt
double zahl = readDouble("Gib eine Zahl ein"); // Mit Prompt
\end{lstlisting}

\subsubsection*{Beispiele:}
\begin{lstlisting}
// Preis einlesen
double preis = readDouble("Preis in Euro");
double mehrwertsteuer = preis * 0.19;
double gesamtpreis = preis + mehrwertsteuer;
System.out.println("Gesamt: " + gesamtpreis + " Euro");

// Durchschnitt berechnen
double note1 = readDouble("Note 1");
double note2 = readDouble("Note 2");
double note3 = readDouble("Note 3");
double durchschnitt = (note1 + note2 + note3) / 3.0;
System.out.println("Durchschnitt: " + durchschnitt);

// Temperatur
double celsius = readDouble("Temperatur in Celsius");
double fahrenheit = celsius * 9.0 / 5.0 + 32.0;
System.out.println(celsius + "°C = " + fahrenheit + "°F");

// Auch Ganzzahlen moeglich
double wert = readDouble("Zahl");  // Eingabe: 5
// wert = 5.0
\end{lstlisting}

\begin{achtung}
	\begin{itemize}
		\item Je nach Computer funktioniert Komma \texttt{,} oder Punkt \texttt{.} als Trennzeichen
		\item Wenn jemand normalen Text eingibt, stürzt das Programm ab
		\item Du kannst auch ganze Zahlen eingeben - sie werden automatisch zu Kommazahlen
		\item Für die meisten Fälle ist \texttt{readDouble()} besser als \texttt{readFloat()}
	\end{itemize}
\end{achtung}

% ============================================================
\newpage


\subsection*{readByte() - Sehr kleine Ganzzahl einlesen}
\addcontentsline{toc}{subsection}{readByte() - Sehr kleine Ganzzahl einlesen}

\subsubsection*{Was macht die Methode?}\\
Liest eine sehr kleine Ganzzahl von der Konsole ein und gibt sie als \texttt{byte} zurück.

\subsubsection*{Syntax:}
\begin{lstlisting}
byte zahl = readByte();                    // Ohne Prompt
byte zahl = readByte("Gib eine Zahl ein"); // Mit Prompt
\end{lstlisting}

\subsubsection*{Beispiele:}
\begin{lstlisting}
// Alter (wenn < 127)
byte alter = readByte("Alter");

// Prozent (0-100)
byte prozent = readByte("Prozent (0-100)");
\end{lstlisting}

\begin{achtung}
	\begin{itemize}
		\item Wertebereich: nur -128 bis 127
		\item Bei zu großen Zahlen: Fehler
		\item Meist wird \texttt{readInt()} bevorzugt
	\end{itemize}
\end{achtung}

% ============================================================
\newpage


\subsection*{readShort() - Kleine Ganzzahl einlesen}
\addcontentsline{toc}{subsection}{readShort() - Kleine Ganzzahl einlesen}

\subsubsection*{Was macht die Methode?}\\
Liest eine kleine Ganzzahl von der Konsole ein und gibt sie als \texttt{short} zurück.

\subsubsection*{Syntax:}
\begin{lstlisting}
short zahl = readShort();                    // Ohne Prompt
short zahl = readShort("Gib eine Zahl ein"); // Mit Prompt
\end{lstlisting}

\subsubsection*{Beispiele:}
\begin{lstlisting}
// Jahr
short jahr = readShort("Geburtsjahr");

// Distanz in Metern
short meter = readShort("Distanz in Metern");
\end{lstlisting}

\begin{achtung}
	\begin{itemize}
		\item Wertebereich: -32.768 bis 32.767
		\item Bei zu großen Zahlen: Fehler
		\item Meist wird \texttt{readInt()} bevorzugt
	\end{itemize}
\end{achtung}

% ============================================================
\newpage


\subsection*{readLong() - Sehr große Ganzzahl einlesen}
\addcontentsline{toc}{subsection}{readLong() - Sehr große Ganzzahl einlesen}

\subsubsection*{Was macht die Methode?}\\
Liest eine sehr große Ganzzahl von der Konsole ein und gibt sie als \texttt{long} zurück.

\subsubsection*{Syntax:}
\begin{lstlisting}
long zahl = readLong();                    // Ohne Prompt
long zahl = readLong("Gib eine Zahl ein"); // Mit Prompt
\end{lstlisting}

\subsubsection*{Beispiele:}
\begin{lstlisting}
// Sehr große Zahlen
long weltbevoelkerung = readLong("Weltbevoelkerung");

// Millisekunden
long millis = readLong("Millisekunden");
\end{lstlisting}

\begin{achtung}
	\begin{itemize}
		\item Für sehr große Zahlen (größer als ca. 2 Milliarden)
		\item Wertebereich: ca. -9,2 $\times$ 10$^{18}$ bis 9,2 $\times$ 10$^{18}$
		\item Für normale Zahlen: \texttt{readInt()} verwenden
	\end{itemize}
\end{achtung}

% ============================================================
\newpage


\subsection*{readFloat() - Einfache Kommazahl einlesen}
\addcontentsline{toc}{subsection}{readFloat() - Einfache Kommazahl einlesen}

\subsubsection*{Was macht die Methode?}\\
Liest eine Kommazahl mit einfacher Genauigkeit von der Konsole ein und gibt sie als \texttt{float} zurück.

\subsubsection*{Syntax:}
\begin{lstlisting}
float zahl = readFloat();                    // Ohne Prompt
float zahl = readFloat("Gib eine Zahl ein"); // Mit Prompt
\end{lstlisting}

\subsubsection*{Beispiele:}
\begin{lstlisting}
// Gewicht
float gewicht = readFloat("Gewicht in kg");

// Groesse
float groesse = readFloat("Groesse in m");
float bmi = gewicht / (groesse * groesse);
\end{lstlisting}

\begin{achtung}
	\begin{itemize}
		\item Weniger genau als \texttt{double} (ca. 6-7 Dezimalstellen)
		\item Meist wird \texttt{readDouble()} bevorzugt
		\item Dezimaltrennzeichen: Komma oder Punkt (je nach System)
	\end{itemize}
\end{achtung}

% ============================================================
\newpage


\subsection*{readIntArray() - Ganzzahl-Array einlesen}
\addcontentsline{toc}{subsection}{readIntArray() - Ganzzahl-Array einlesen}

\subsubsection*{Was macht die Methode?}\\
Liest mehrere Ganzzahlen auf einmal ein und gibt sie als Array zurück.
Die Zahlen müssen mit Leerzeichen getrennt werden.

\subsubsection*{Syntax:}
\begin{lstlisting}
int[] zahlen = readIntArray();                    // Ohne Prompt
int[] zahlen = readIntArray("Gib Zahlen ein");    // Mit Prompt
\end{lstlisting}

\subsubsection*{Beispiele:}
\begin{lstlisting}
// Mehrere Zahlen auf einmal einlesen
int[] noten = readIntArray("Noten (durch Leerzeichen getrennt)");
// Eingabe: 1 2 3 2 1
// noten = {1, 2, 3, 2, 1}

// Durchschnitt berechnen
int summe = 0;
for (int note : noten) {
    summe += note;
}
double durchschnitt = (double) summe / noten.length;
System.out.println("Durchschnitt: " + durchschnitt);

// Array direkt ausgeben
int[] zahlen = readIntArray("Zahlen");
System.out.println("Anzahl: " + zahlen.length);
for (int i = 0; i < zahlen.length; i++) {
    System.out.println("Zahl " + (i + 1) + ": " + zahlen[i]);
}
\end{lstlisting}

\begin{achtung}
	\begin{itemize}
		\item Die Zahlen müssen mit Leerzeichen getrennt sein (z.B. "1 2 3")
		\item Die Größe des Arrays passt sich automatisch an die Anzahl der Zahlen an
		\item Wenn jemand Text eingibt, stürzt das Programm ab
		\item Alle Zahlen müssen in einer Zeile stehen (Enter beendet die Eingabe)
	\end{itemize}
\end{achtung}

% ============================================================
\newpage


\subsection*{readDoubleArray() - Kommazahl-Array einlesen}
\addcontentsline{toc}{subsection}{readDoubleArray() - Kommazahl-Array einlesen}

\subsubsection*{Was macht die Methode?}\\
Liest mehrere Kommazahlen von der Konsole ein und gibt sie als \texttt{double[]}-Array zurück.
Die Zahlen müssen durch Leerzeichen getrennt eingegeben werden.

\subsubsection*{Syntax:}
\begin{lstlisting}
double[] zahlen = readDoubleArray();                    // Ohne Prompt
double[] zahlen = readDoubleArray("Gib Zahlen ein");    // Mit Prompt
\end{lstlisting}

\subsubsection*{Beispiele:}
\begin{lstlisting}
// Preise einlesen
double[] preise = readDoubleArray("Preise");
// Eingabe: 9.99 14.50 7.25
// preise = {9.99, 14.50, 7.25}

// Gesamtpreis berechnen
double gesamt = 0.0;
for (double preis : preise) {
    gesamt += preis;
}
System.out.println("Gesamt: " + gesamt + " Euro");

// Temperaturwerte
double[] temps = readDoubleArray("Temperaturen");
double min = temps[0];
double max = temps[0];
for (double temp : temps) {
    if (temp < min) min = temp;
    if (temp > max) max = temp;
}
System.out.println("Min: " + min + ", Max: " + max);
\end{lstlisting}

\begin{achtung}
	\begin{itemize}
		\item Zahlen müssen durch \textbf{Leerzeichen} getrennt werden
		\item Dezimaltrennzeichen: Komma oder Punkt (je nach System)
		\item Array-Größe ergibt sich automatisch aus der Anzahl der Eingaben
		\item Bei ungültiger Eingabe: Programm bricht ab
	\end{itemize}
\end{achtung}

% ============================================================
\newpage


\subsection*{readStringArray() - String-Array einlesen}
\addcontentsline{toc}{subsection}{readStringArray() - String-Array einlesen}

\subsubsection*{Was macht die Methode?}\\
Liest mehrere Wörter von der Konsole ein und gibt sie als \texttt{String[]}-Array zurück.
Die Wörter müssen durch Leerzeichen getrennt eingegeben werden.

\subsubsection*{Syntax:}
\begin{lstlisting}
String[] woerter = readStringArray();                    // Ohne Prompt
String[] woerter = readStringArray("Gib Woerter ein");   // Mit Prompt
\end{lstlisting}

\subsubsection*{Beispiele:}
\begin{lstlisting}
// Namen einlesen
String[] namen = readStringArray("Namen");
// Eingabe: Anna Bob Clara
// namen = {"Anna", "Bob", "Clara"}

// Alle Namen ausgeben
System.out.println("Teilnehmer:");
for (String name : namen) {
    System.out.println("- " + name);
}

// Woerter zaehlen
String[] woerter = readStringArray("Satz");
System.out.println("Anzahl Woerter: " + woerter.length);

// Laengstes Wort finden
String[] text = readStringArray("Text");
String laengstes = text[0];
for (String wort : text) {
    if (wort.length() > laengstes.length()) {
        laengstes = wort;
    }
}
System.out.println("Laengstes Wort: " + laengstes);
\end{lstlisting}

\begin{achtung}
	\begin{itemize}
		\item Die Wörter müssen mit Leerzeichen getrennt sein
		\item Wenn du ganze Sätze einlesen willst, benutze lieber \texttt{readString()}
		\item Die Größe des Arrays passt sich automatisch an die Anzahl der Wörter an
		\item Für Sätze: \texttt{readString()} benutzen und dann mit \texttt{.split(" ")} aufteilen
	\end{itemize}
\end{achtung}

% ============================================================
\newpage


\subsection*{Weitere Array-Methoden}
\addcontentsline{toc}{subsection}{Weitere Array-Methoden}

Der ConsoleReader bietet auch Methoden für andere Datentypen:

\begin{itemize}
	\item \texttt{readByteArray()} - Liest \texttt{byte[]}-Array ein
	\item \texttt{readShortArray()} - Liest \texttt{short[]}-Array ein
	\item \texttt{readLongArray()} - Liest \texttt{long[]}-Array ein
	\item \texttt{readFloatArray()} - Liest \texttt{float[]}-Array ein
\end{itemize}

Diese funktionieren analog zu \texttt{readIntArray()} und \texttt{readDoubleArray()}.

\subsubsection*{Beispiel:}
\begin{lstlisting}
// Byte-Array
byte[] bytes = readByteArray("Bytes");

// Long-Array fuer grosse Zahlen
long[] ids = readLongArray("IDs");

// Float-Array
float[] gewichte = readFloatArray("Gewichte");
\end{lstlisting}

% ============================================================
\newpage


\section*{Bildverarbeitung}
\addcontentsline{toc}{section}{Bildverarbeitung}

Die \texttt{Picture}-Klasse ist ein Werkzeug zum Arbeiten mit Bildern.
Bilder werden als zweidimensionale Arrays gespeichert, wobei jeder Wert eine Farbe darstellt.
Um sie zu benutzen, schreibst du am Anfang:

\begin{lstlisting}
import static prog.Picture.*;
\end{lstlisting}

% ============================================================
\newpage


\subsection*{show() - Bild anzeigen}
\addcontentsline{toc}{subsection}{show() - Bild anzeigen}

\subsubsection*{Was macht die Methode?}\\
Zeigt ein Bild (als \texttt{int[][]}-Array) in einem Fenster an.

\subsubsection*{Syntax:}
\begin{lstlisting}
show(bildArray);                    // Ohne Fenstertitel
show(bildArray, "Fenstertitel");    // Mit Fenstertitel
\end{lstlisting}

\subsubsection*{Beispiele:}
\begin{lstlisting}
// Bild aus Datei laden und anzeigen
int[][] bild = loadFromFileSystem("pfad/zum/bild.jpg");
show(bild);

// Mit eigenem Fenstertitel
int[][] foto = loadFromFileSystem("foto.png");
show(foto, "Mein Urlaubsfoto");

// Mehrere Bilder anzeigen
int[][] bild1 = loadResource("bild1.jpg");
int[][] bild2 = loadResource("bild2.jpg");
show(bild1, "Original");
show(bild2, "Bearbeitet");

// Selbst erstelltes Bild anzeigen
int[][] neuBild = new int[100][100];
for (int y = 0; y < neuBild.length; y++) {
    for (int x = 0; x < neuBild[0].length; x++) {
        neuBild[y][x] = color(255, 0, 0);  // Rot
    }
}
show(neuBild, "Rotes Quadrat");
\end{lstlisting}

\begin{achtung}
	\begin{itemize}
		\item Das Array hat das Format \texttt{int[Höhe][Breite]} - erst Höhe, dann Breite
		\item Jede Zahl im Array stellt eine Farbe dar
		\item Das Fenster bleibt offen, bis du es schließt
		\item Wenn du keinen Titel angibst, wird ein Standardtitel verwendet
	\end{itemize}
\end{achtung}

% ============================================================
\newpage


\subsection*{loadFromFileSystem() - Bild von Festplatte laden}
\addcontentsline{toc}{subsection}{loadFromFileSystem() - Bild von Festplatte laden}

\subsubsection*{Was macht die Methode?}\\
Lädt ein Bild von der Festplatte und gibt es als \texttt{int[][]}-Array zurück.

\subsubsection*{Syntax:}
\begin{lstlisting}
int[][] bild = loadFromFileSystem("pfad/zur/datei.jpg");
\end{lstlisting}

\subsubsection*{Beispiele:}
\begin{lstlisting}
// Absoluter Pfad
int[][] bild1 = loadFromFileSystem("/home/user/bilder/foto.jpg");
show(bild1);

// Relativer Pfad (relativ zum Projektverzeichnis)
int[][] bild2 = loadFromFileSystem("bilder/landschaft.png");
show(bild2, "Landschaft");

// Bild laden und Informationen ausgeben
int[][] foto = loadFromFileSystem("foto.jpg");
int hoehe = foto.length;
int breite = foto[0].length;
System.out.println("Bildgroesse: " + breite + " x " + hoehe);

// Bild laden und verarbeiten
int[][] original = loadFromFileSystem("bild.jpg");
int[][] bearbeitet = new int[original.length][original[0].length];
// ... Bildverarbeitung ...
show(original, "Original");
show(bearbeitet, "Bearbeitet");
\end{lstlisting}

\begin{achtung}
	\begin{itemize}
		\item Funktioniert mit den gängigen Bildformaten: JPG, PNG, BMP, GIF
		\item Du kannst den vollständigen Pfad oder einen relativen Pfad angeben
		\item Wenn die Datei nicht existiert, stürzt das Programm ab
		\item Für Bilder, die zu deinem Programm gehören: benutze \texttt{loadResource()}
	\end{itemize}
\end{achtung}

% ============================================================
\newpage


\subsection*{loadResource() - Bild aus Projekt laden}
\addcontentsline{toc}{subsection}{loadResource() - Bild aus Projekt laden}

\subsubsection*{Was macht die Methode?}\\
Lädt ein Bild, das zu deinem Projekt gehört, und gibt es als Array zurück.
Das Bild muss im selben Ordner liegen wie dein Programm.

\subsubsection*{Syntax:}
\begin{lstlisting}
int[][] bild = loadResource("dateiname.jpg");
\end{lstlisting}

\subsubsection*{Beispiele:}
\begin{lstlisting}
// Bild direkt aus dem Package laden
int[][] logo = loadResource("logo.png");
show(logo, "Logo");

// Mehrere Bilder laden
int[][] bild1 = loadResource("bild1.jpg");
int[][] bild2 = loadResource("bild2.jpg");
show(bild1, "Bild 1");
show(bild2, "Bild 2");

// Bild laden und Groesse pruefen
int[][] foto = loadResource("foto.png");
if (foto.length > 0 && foto[0].length > 0) {
    System.out.println("Bild erfolgreich geladen");
    show(foto);
}

// Typische Verwendung in einfachen Projekten
// Dateistruktur:
// src/
//   prog/
//     Main.java
//     bild.jpg  <- Bild liegt direkt neben der Klasse

// Variante 1: Relativ
int[][] bild = loadResource("bild.jpg");
show(bild);

// Variante 2: Absolut vom Classpath-Root
int[][] bild2 = loadResource("/prog/bild.jpg");
show(bild2);
\end{lstlisting}

\begin{achtung}
	\begin{itemize}
		\item Nur für Bilder benutzen, die zu deinem Projekt gehören
		\item Das Bild muss im selben Ordner liegen wie dein Programm
		\item Gib nur den Dateinamen an, keinen Pfad
		\item Für Bilder von woanders auf deinem Computer: benutze \texttt{loadFromFileSystem()}
	\end{itemize}
\end{achtung}

% ============================================================
\newpage


\subsection*{color() - Farbwert erstellen}
\addcontentsline{toc}{subsection}{color() - Farbwert erstellen}

\subsubsection*{Was macht die Methode?}\\
Erstellt einen RGB-Farbwert aus Rot-, Grün- und Blau-Anteilen.

\subsubsection*{Syntax:}
\begin{lstlisting}
int farbwert = color(rot, gruen, blau);  // Werte: 0-255
\end{lstlisting}

\subsubsection*{Beispiele:}
\begin{lstlisting}
// Grundfarben
int rot = color(255, 0, 0);
int gruen = color(0, 255, 0);
int blau = color(0, 0, 255);

// Mischfarben
int gelb = color(255, 255, 0);      // Rot + Gruen
int magenta = color(255, 0, 255);   // Rot + Blau
int cyan = color(0, 255, 255);      // Gruen + Blau

// Graustufen
int schwarz = color(0, 0, 0);
int weiss = color(255, 255, 255);
int grau = color(128, 128, 128);

// Bild mit Farbe fuellen
int[][] bild = new int[100][200];
int orange = color(255, 165, 0);
for (int y = 0; y < bild.length; y++) {
    for (int x = 0; x < bild[0].length; x++) {
        bild[y][x] = orange;
    }
}
show(bild, "Orange");

// Farbverlauf erstellen
int[][] verlauf = new int[256][100];
for (int y = 0; y < verlauf.length; y++) {
    int farbwert = color(y, 0, 255 - y);  // Von Blau zu Rot
    for (int x = 0; x < verlauf[0].length; x++) {
        verlauf[y][x] = farbwert;
    }
}
show(verlauf, "Farbverlauf");

// Pixel-Farbe setzen
int[][] foto = loadResource("foto.jpg");
foto[50][100] = color(255, 0, 0);  // Roten Punkt setzen
show(foto);
\end{lstlisting}

\begin{achtung}
	\begin{itemize}
		\item Jede Farbe (Rot, Grün, Blau) kann einen Wert von 0 bis 255 haben
		\item 0 bedeutet "gar nichts von dieser Farbe", 255 bedeutet "maximale Helligkeit"
		\item RGB steht für: Rot, Grün, Blau
		\item Die Methode gibt eine einzige Zahl zurück, die alle drei Farben enthält
	\end{itemize}
\end{achtung}

% ============================================================
\newpage


\subsection*{red(), green(), blue() - Farbkanäle extrahieren}
\addcontentsline{toc}{subsection}{red(), green(), blue() - Farbkanäle extrahieren}

\subsubsection*{Was machen die Methoden?}\\
Holen den Rot-, Grün- oder Blau-Anteil aus einer Farbe heraus.

\subsubsection*{Syntax:}
\begin{lstlisting}
int rotWert = red(farbwert);      // Rot-Anteil (0-255)
int gruenWert = green(farbwert);  // Gruen-Anteil (0-255)
int blauWert = blue(farbwert);    // Blau-Anteil (0-255)
\end{lstlisting}

\subsubsection*{Beispiele:}
\begin{lstlisting}
// Farbkanäle eines Farbwerts extrahieren
int orange = color(255, 165, 0);
int r = red(orange);      // r = 255
int g = green(orange);    // g = 165
int b = blue(orange);     // b = 0
System.out.println("RGB: " + r + ", " + g + ", " + b);

// Pixel-Farbe analysieren
int[][] bild = loadResource("foto.jpg");
int pixel = bild[50][100];  // Pixel an Position (100, 50)
System.out.println("Rot: " + red(pixel));
System.out.println("Gruen: " + green(pixel));
System.out.println("Blau: " + blue(pixel));

// Bild in Graustufen umwandeln
int[][] original = loadResource("bild.jpg");
int[][] graustufen = new int[original.length][original[0].length];
for (int y = 0; y < original.length; y++) {
    for (int x = 0; x < original[0].length; x++) {
        int pixel = original[y][x];
        int r = red(pixel);
        int g = green(pixel);
        int b = blue(pixel);
        int grau = (r + g + b) / 3;  // Durchschnitt
        graustufen[y][x] = color(grau, grau, grau);
    }
}
show(original, "Original");
show(graustufen, "Graustufen");

// Nur Rot-Kanal behalten
int[][] nurRot = new int[original.length][original[0].length];
for (int y = 0; y < original.length; y++) {
    for (int x = 0; x < original[0].length; x++) {
        int r = red(original[y][x]);
        nurRot[y][x] = color(r, 0, 0);
    }
}
show(nurRot, "Nur Rot-Kanal");

// Farbkanäle vertauschen (Rot <-> Blau)
int[][] vertauscht = new int[original.length][original[0].length];
for (int y = 0; y < original.length; y++) {
    for (int x = 0; x < original[0].length; x++) {
        int pixel = original[y][x];
        int r = red(pixel);
        int g = green(pixel);
        int b = blue(pixel);
        vertauscht[y][x] = color(b, g, r);  // Blau und Rot tauschen
    }
}
show(vertauscht, "Vertauschte Kanäle");

// Helligkeit eines Pixels berechnen
int pixel = bild[0][0];
int helligkeit = (red(pixel) + green(pixel) + blue(pixel)) / 3;
System.out.println("Helligkeit: " + helligkeit);
\end{lstlisting}

\begin{achtung}
	\begin{itemize}
		\item Du bekommst immer eine Zahl zwischen 0 und 255 zurück
		\item Sehr nützlich, wenn du Bilder bearbeiten willst
		\item Zusammen mit \texttt{color()} kannst du einzelne Farbanteile ändern
		\item Häufige Anwendung: Bilder in Graustufen umwandeln oder Farbfilter erstellen
	\end{itemize}
\end{achtung}

% ============================================================
\newpage


\section*{Häufige Fehler bei Strings}
\addcontentsline{toc}{section}{Häufige Fehler bei Strings}

\subsection*{String.length ohne Klammern}

\begin{lstlisting}
// FALSCH
String text = "Hallo";
int laenge = text.length;  // Fehler: Strings brauchen ()

// RICHTIG
int laenge = text.length();    // String MIT Klammern
\end{lstlisting}

\subsection*{Index-Grenzen bei charAt}

\begin{lstlisting}
String wort = "Test";  // length = 4, Indizes: 0, 1, 2, 3
char c = wort.charAt(4);  // Fehler! Index 4 existiert nicht

// RICHTIG
char letztes = wort.charAt(wort.length() - 1);  // Index 3
char erstes = wort.charAt(0);                   // Index 0
\end{lstlisting}

\subsection*{char vs. String}

\begin{lstlisting}
String text = "A";
char buchstabe = 'A';

// FALSCH
if (text.charAt(0) == "A") { }  // Vergleicht char mit String

// RICHTIG
if (text.charAt(0) == 'A') { }  // char mit char vergleichen
\end{lstlisting}
